\section*{Step 34: Instance-Uniqueness Up To Equivalence}

\paragraph{Data.}
Let \(q_{\rm QM}=(d_0,d_1,d_2)\), \(q_{\rm GR}=(d_0,d_2,d_3)\), and
\[
  L=(d_0,d_1,d_2,d_3).
\]
All maps are finite coordinate quotients on the declared carrier.

\paragraph{Join.}
Both endpoint quotients factor through \(L\) with residual \(0\).  Since the row directions of \(L\) are exactly the union of endpoint-distinguished directions,
\[
  \{d_0,d_1,d_2\}\cup\{d_0,d_2,d_3\}=\{d_0,d_1,d_2,d_3\},
\]
\(L\) is the coordinate-lattice join of \(q_{\rm QM}\) and \(q_{\rm GR}\).

\paragraph{Minimality.}
Dropping any one of \(d_0,d_1,d_2,d_3\) leaves at least one endpoint quotient with positive linear factorization residual.  Thus no proper coordinate subobject of \(L\) is a common refinement.

\paragraph{Non-minimal refinements.}
The five-mode object \(M=L+d_4\) is admissible, but \(L\) factors through \(M\) while \(M\) does not factor through \(L\).  Hence \(M\) is a non-minimal refinement rather than a second minimal mediator.

\paragraph{Union and equivalence.}
The direct-sum union has rank \(4\) in dimension \(6\), so it fails the duplicate-coordinate admissibility check.  A relabeled copy \(L'=PL\) is isomorphic to \(L\), with residual \(0\) in both directions under the exhibited permutation.

\paragraph{Conclusion.}
Within the declared finite coordinate-lattice grammar, \(L\) is the unique minimal admissible common refinement of \(q_{\rm QM}\) and \(q_{\rm GR}\), up to isomorphism.  This is not a claim that \(L\) is the only admissible refinement; it is the unique minimal one in this grammar.
